%% ma: L11-B11-0002
%% lop: 11
%% chuong: 4
%% bai: 11
%% dang: 11.11.2
%% loai: DS
%% muc: [NB, TH, VD, TH]
%% dap_an: "ĐSĐĐ"
%% nguon: "de-cuong-giua-ky-1-toan-11-nam-2026-2027-truong-thpt-hoang-van-thu-ha-noi-2-8.pdf (M08, Đề cương giữa kì 1 Toán 11 THPT Hoàng Văn Thụ, Hà Nội), Phần 2 Câu 8"
%% kien_thuc: "Bài 10 (giao tuyến, giao điểm), Bài 11 (hai đường thẳng song song, định lí về giao tuyến)"
%% trang_thai: chinh-thuc
%% ghi_chu: "Đã duyệt 29/10/2026."
\begin{ex}
Cho hình chóp $S.ABCD$ có đáy $ABCD$ là một hình bình hành tâm $O$. Gọi $I$, $K$ lần lượt là trung điểm của $SB$ và $SD$. Khi đó:
\choiceTF{\True $SO$ là giao tuyến của $(SAC)$ và $(SBD)$}{Giao điểm $J$ của $SA$ với $(CKB)$ thuộc đường thẳng đi qua $K$ và song song với $DC$}{\True Giao tuyến của $(OIA)$ và $(SCD)$ là đường thẳng đi qua $C$ và song song với $SD$}{\True $CD\parallel IJ$}
\loigiai{a) Đúng: $S$, $O$ là hai điểm chung.

b) Sai: $(CKB)$ chứa $BC\parallel AD$, $(SAD)$ chứa $AD$ và có điểm chung $K$, nên $(CKB)\cap(SAD)$ là đường thẳng qua $K$ song song với $AD$; đường này cắt $SA$ tại $J$. Vậy $J$ thuộc đường thẳng qua $K$ song song với $AD$, không phải $DC$.

c) Đúng: $A$, $O$, $C$ thẳng hàng nên $C\in(OIA)$; $OI$ là đường trung bình của $\triangle SBD$ nên $OI\parallel SD$. Hai mặt phẳng $(OIA)$, $(SCD)$ chứa hai đường thẳng song song $OI$, $SD$ và có điểm chung $C$, nên giao tuyến qua $C$ song song với $SD$.

d) Đúng: $KJ\parallel AD$ và $K$ là trung điểm $SD$ nên $J$ là trung điểm $SA$; $IJ$ là đường trung bình của $\triangle SAB$ nên $IJ\parallel AB\parallel CD$.}
\end{ex}
